% !TeX root = ../papers/nc-p3.tex
\section{二本の有理正規三次からなる六次点グラフの全排除}
\label{sec:binary-sextic-exclusion}
\begin{theorem}\label{thm:binary-sextic-exclusion}
四辺が点ファイバーの標準AS窓の共通点グラフは、二本の有理正規三次 $C_1,C_2$ の和 $Y$ にはならない。
四点が重なる場合も含める。
\end{theorem}
\begin{proof}[曲線の切断系と可動な三割線]
純性と種数公式より $T=C_1\cap C_2$ は長さ4である。$T$ は両射影直線上のCartier因子であり、
$Y$ はその貼合せである。各共通点の形式環は二つの滑らかな枝を同じ長さ $m$ の点で貼った
$k[[x,y]]/(y(y-x^m))$ で、$Y$ はGorenstein、$\omega_Y|_{C_j}=\mathcal O(2)$。
各 $L_i$ の両成分上の三次切断は $T$ 上の四切断に同型に制限される。
従って四文字は $L_i$ の完全切断系であり、$H^1(L_i)=0$。

$T$ の貼合せが両直線間の射影同型に延びれば、$Y$ は直線への有限平坦二重被覆となる。
局所的には同じ直線の函数二つを $T$ 上で一致させる代数であり、trace零部分は $\mathcal O(-4)$ である。
被覆鉛筆を $J$ として $\omega_Y=J^2$。$L_i/J^2$ の切断は片側で零ならもう一方の
次数1の切断も $T$ 上で消えて零になる。非零なら長さ2のCartier零因子を持ち、
$L_i$ の三切断がその上で消え、埋込みに反する。従って $f_*L_i=\mathcal O(1)^2$。
多重次数零の束の非零切断は連結な被約曲線上で生成元なので、
定理\ref{thm:smooth-sextic-exclusion}の有限平坦二重被覆の計算で三次元数は高々18となり不可能である。
この計算で二鉛筆が隣接して等しいときは、Segre二次曲面の対角が入る。
実際二曲線の和を含む二次式が二本あれば、独立二式の次数4または共通平面因子と非平面な台の矛盾になる。

以後貼合せが射影同型に延びないとする。各位置の $Y$ は二次曲面に入らない。
二平面では非平面な成分を含めず、錐では純六次因子の反射的イデアルは
$\mathcal O_Q(-3)$ で算術種数4になる。滑らかな二次曲面で各三次は $(1,2)$ または $(2,1)$。
両型が異なれば算術種数4、同じなら一方の射影が各成分で同型になり、上の貼合せが延びてしまう。
標準系は各成分の二次完全系に同型に制限される。実際、片方で零の標準微分は他方で正則な微分になるからである。
二つの標準像は異なる滑らかな平面二次曲線となる。同じなら両成分で同型な標準射影が貼合せを延ばす。
相異なるなら像の平面四次と $Y$ の算術種数が一致するので有限双有理な標準写像は同型である。

$\tau=\sigma^{-1}$ と書く。位置0で一般の滑らかな $p$ を固定する。
反対二点縮約の階数1以下の因子 $D_p$ は $Y$ 上で
\[
                   \mathcal O(D_p)=K^{-1}(p+\tau p)
\tag{*}
\]
である。定理\ref{thm:opposite-curve-divisor}の二つの核直線の共通零は、
固定した $p$ の場合には滑らかな二点 $p,\tau p$（同一点でもよい）の可逆イデアルである。
各成分で一般階数は2である。階数1の外積なら、両端の三つずつの函数の積は少なくとも五つを張るが、
変数側は各成分が張る四文字だけだからである。逆に $q$ を固定した場合にもこの議論が通るので、
一般の $p$ では全ての交点を避けて因子が得られる。
二成分上で二点とも動かす際、対角や $\tau$ のグラフがある成分の積ではその因子を核直線から除き、
異なる成分の像の座標が一致する有限個の点は除いて行列式を取る。
各射影直線の積は滑らかなので、二次小行列と核の外積の比は余次元2を越えて延びる
（極は余次元1で検出される）。従って階数低下の因子の二次数は、
当該成分積に存在する対角と $\tau$ のグラフの二次数の和である。
どの縦横ファイバーも因子として入らない。実際、交点を除いたその一般点では
既述の三次元函数空間の積の下界が使える。従って一方の射影に次数1なら他方にも次数1で射影同型のグラフとなる。

$D_p$ は位置2で $Y$ の三割線に乗ることを示す。
反対縮約の $3\times3$ 行列式 $S_p$ は一般の $p$ で非零である。
どの $R_i$ にも階数1関係はない（一本の三次だけでも全位置で空間を張る）。
従って縮約の中支持は3または4であり、中支持4の場合は定理\ref{thm:finite-cubic-separators}の
UFDでの余因子行列の証明を使える。中支持3で三次式が恒等零なら、
定理\ref{thm:ternary-point-subalgebras}の三支持の三文字部分代数は平面グラフを持つ
（階数1縮約がなく、その中央の行列式が恒等零）。従って第一と中央の三文字での平面射影が
$Y$ 全体で共通になる。その第一の中心は $j_1(p)$ で、中央の中心を $b$ とする。
$p$ の属する三次の射影像は二次曲線だから、$b$ も同じ成分に属する。
外からの三次の射影が非退化な二次曲線になることは次数から不可能である。
同じ比較を他方で行えば $b\notin T$ でもある。
全 $Y$ で基点を除いた偏極は従って $L_1(-p)=L_2(-b)$ である。
二つの一般の $p$ で比較すると、その差は一成分上の二次因子の線型同値となる。
他成分で定数の有理函数が長さ4の $T$ でその値を取り、極は二次以下なので定数である。
従って因子として等しく、$p$ が動くことから $b=p$、$L_1=L_2$ となる。
二次曲面がないので全関係が交換子になり矛盾する。

$Y$ は四つの三次最大小行列で切られる。実際二次式がなく
$h^0(L^2)=10$（両側の六次切断は長さ4へ全射）、
$H^1(L^n)=0$ ($n>0$) なので、基点自由鉛筆の完全列から全高次の制限も全射である。
従って定理\ref{thm:nonhyperelliptic-sextic-exclusion}のHilbert--Burchの証明を使える。
一般の $D_p$ は滑らかな部分にある。$S_p$ はこの三次系に入り $D_p$ で特異なので、
同定理の $3\times4$ 線型行列での核の議論により $D_p$ の張る直線上で
四つの最大小行列は一本の三次式の倍数となる。$Y$ 自体は直線を含まず、交わりはちょうど長さ3である。

\smallskip
\noindent\textit{四移動の位数。}
$\tau$ が一成分を点ごとに固定すれば、標準平面の滑らかな二次曲線を点ごとに固定するので全体で恒等である。
非恒等の場合は定理\ref{thm:fixed-differential}により固定スキームは長さ4のCartier因子である。
従って標準平面の射影変換は固定直線を持つ（孤立固定スキームの長さは3）。
非半単純な固定直線型は推移変換であり、不変な四次式は二変数だけになって二次曲線二本の積にならない。
半単純の場合、固定直線の切断が既に長さ4であるため中心は曲線外で、固定因子はこの標準切断であり $K^2=1$。
固定直線を $u=0$ とすれば四次式に $u^4$ と $u$ を含まない部分が共にあり、倍率の位数は2または4である。
位数4は定理\ref{thm:integral-sextic-exclusion}の巡回固有束の議論で除ける。
詳しくは標準式が $u^4=F_4(v,w)$ となり、$d_i=L_{i+1}/L_i$ は不変でノルムにより $d_i^4=1$、
従って $\tau^*L_i=L_i$。ノルムはここでも有限平坦な巡回四重射の一般軌道上の積で、
Cartier因子のノルム・引戻しの式は滑らかな点を台とする因子で確認できる
（十分正の束を用い特異点を避けた二因子の差で任意の線束を表せる）。
$P=L_2/\omega_Y$ は多重次数 $(1,1)$ で $H^0(P)=H^1(P)=0$。
非零切断は他成分で恒等零なら長さ4で消えるので零、さもなければ長さ2のCartier零因子で埋込みに反する。
よって不変な $P$ を線型化した順像の四固有束は $\mathcal O(-1)$ であり、これは一般軌道が自由推移的なので可約でも同じである。
$u$ の四つの一次乗法の積が $F_4$ で、根の上で $k[u]/u^4$ の自由階数1加群になるためどの二因子も根を共有しない。
すると全根が単純で標準曲線は滑らかになり矛盾する。以後 $\tau$ は恒等または対合である。

\smallskip
\noindent\textit{成分を保つ場合。}
一般の $p\in C_1$ に対して $(*)$ の多重次数は $(2,0)$。
三割線の残る点 $r$ は必ず $C_2\setminus T$ にある。有理正規三次 $C_1$ と直線は長さ3では交われず、
共通点が三番目となることもない。従って $D_p+r$ は滑らかな台上のCartier因子である。
$L_2(-D_p-r)$ には直線を含む平面から二切断があるのでRiemann--Rochにより
$D_p+r\sim P+s$ となる滑らかな $s\in C_1$ が得られる。実際残余の束は多重次数 $(1,0)$ で、
非零切断は $C_2$ で恒等零になれず生成元となり、$C_1$ 上の一点だけで消える。
$P$ は非有効で双対残余も非有効だから $|P+s|$ は一元で、$D_p$ が動く以上 $s$ は非定数である。
$(*)$ より $p+\tau p+r\sim PK+s$ でもある。

多重次数 $(1,1)$ の束 $P'$ について、$|P'+s|$ の因子の $C_2$ 上の一点がどう動くかを記す。
両正規化上の一次切断の $T$ での積
$H^0(\mathcal O_{C_1}(1))\otimes H^0(\mathcal O_{C_2}(1))\to H^0(\mathcal O_T(1,1))$ は同型である。
核があれば $(1,1)$ 形式が $T$ で消えるが、階数2なら貼合せが射影同型に延び、
階数1なら二因子で高々長さ2しか消えないからである。
$C_1$ の二次切断の $T$ 上への像は超平面で、その零化汎関数 $\lambda$ は各局所socleで非零である。
実際長さ3部分では二次切断の評価が同型である。従って $\lambda$ は $T$ 上で完全積対合を与える。
各束の自明化を選べば、$P'$ の貼合せは長さ4の環の単元 $v$ で、全体倍率を除いて決まる。
$s$ で消える一次式 $l_s$ と $C_2$ の一次切断 $m$ に対し、$m$ が $|P'+s|$ の切断の制限となる条件は
$\lambda(v l_s m)=0$ である。$C_1$ の二次切断にその値が持ち上がる条件そのものである。
上の二つの同型から $v\mapsto (l,m\mapsto\lambda(vlm))$ は四次元空間の線型同型である。
$P'=P$ では行列は可逆である。階数が落ちれば $s$ に関係なく条件を満たす $m\ne0$ があり、
$vm$ は $C_1$ の一次切断の評価像に入る。実際その像は両方の $l$ との積を $\lambda$ が消す二次元空間である。
これは非有効性に反する。従って $r$ は $s$ の射影同型像である。
$P'=PK$ も同じ非定数の対応を全ての一般の $s$ で与えるので、その行列は比例し、$P=PK$、$K=1$ である。
$(*)$ の二次因子は $D_p=p+\tau p$ として一致する。他方の成分で定数となる有理函数の極次数2と
$T$ の長さ4を再び比較すればよい。同じ結論が両成分で成立する。

$\tau$ が非自明な対合なら両成分上で非自明であり、
$C_1$ の $p,\tau p$ を結ぶ割線の和は滑らかな二次曲面となる。
実際、座標を $t\mapsto-t$ とすれば割線は $[1:0:t^2:0],[0:1:0:t^2]$ の結ぶ直線である。
各割線が $C_2$ にも会い、その交点が動く（一定なら一本の母線を除く全てにも含まれてしまう）。
よって $Y$ が二次曲面に入り矛盾する。

$\tau=1$ ならば両三次の各接線は相手に交わる。三次の接線面を二項三次の
$[u^2v]$ ($u,v$ は二項一次式) で表示すると、対角が元の三次で、
その外では重根 $u$ と残りの根 $v$ は一意である。
従って $C_2$ は $u\mapsto[u^2 h(u)]$ と書ける。実際、$C_2$ から両射影直線への二つの射の次数を
$e,f$ とすると $2e+f=3$ である。接点は動くので $e\ge1$、ゆえに $e=f=1$ で $h$ は射影同型となる。
$h=1$ なら同じ三次であり、$h$ の固有点だけが二曲線の共通点になる。
共通点が一つなら $\operatorname{Pic}^0Y$ は長さ4の局所環の単元を定数倍で割った群なので
標数零で捩れを持たない。既に $K=1$、四移動が恒等なので $L_{i+1}/L_i$ の四乗は1、全偏極が等しく不可能である。
共通点が二つなら $h$ は半単純である。
$C_2=H C_1$ とすると $H$ は対角行列で、その四成分は等差数列
$1,1+d,1+2d,1+3d$ に規格化できる。これは $u^2h(u)$ の二項展開による
（$\operatorname{Sym}^3$ の単項式基底を両側で揃える）。全て非零であり $d\ne0$。
逆に $C_1$ も $C_2$ の接線面に入るため $H^{-1} C_1$ は $C_1$ の接線面に入る。
対角行列 $H^{-1}$ は二つの共通端点を別々に固定するので、重根のパラメータ変更も残りの根の対応もその二点を別々に固定する。
従って $v_j=(1+jd)^{-1}$ は他の等差数列に等比数列を掛けたものである。
ある $b\ne0$ で $v_{j+2}-2b v_{j+1}+b^2v_j=0$ ($j=0,1$) が必要だが、
分母を除き $b$ を消去した終結式は $8d^6(1+d)(1+2d)$ で、矛盾する。

\smallskip
\noindent\textit{成分を交換する対合。}
標準平面の同次座標で $\tau:[u:v:w]\mapsto[-u:v:w]$ とし、一方の二次曲線を
$u^2+u l(v,w)+q(v,w)=0$ とする。もう一方は $u$ の符号を替えた式であり、$l\ne0$。
$T$ は各二次曲線上で $T_0+T_1$、ここで $T_0=(u=0)$、$T_1=(l=0)$ はともに次数2である。
二因子は相異なる。実際同じ切断因子を与える平面線型式なら二次曲線上で比例する。
$T$ を $C_1$ 上で見る埋込みを $x$、$C_2$ 側の埋込みを $\tau^{-1}$ で $C_1$ に運んだものを $y$ とする。
$x=y$ のイデアルは $T$ 上で $(u)$ で、その零化イデアルは $T_1$ のイデアル $(l)$ である。
中心外で射影座標の差を取れば $(u)$ となり、$C_1$ 上で $T$ は $ul$ のCartier因子だからである。
$T_1$ の二点が異なれば $y$ は $x$ の二点を交換する。
二重点なら $u=0=l$ の点で接線 $l=0$ が中心を通り、その接方向に対合は倍率 $-1$ を誘導する。
この場合 $y$ は $x$ に対し $T_1$ の一階jetの符号を替える。

前述の二次数より $p\in C_1$ に対して $D_p=A(p)+B(p)$ で、$A:C_1\to C_1$、$B:C_1\to C_2$ は射影同型である。
$(*)$ の線束同値を $T$ 上の貼合せで比較すると、
\[
         A^{-1}x+y=x+B^{-1}\tau y\qquad\text{in }\operatorname{Sym}^2(C_1)(T).
\tag{**}
\]
詳細には $[x,A(p)][y,p]$ と $[x,p][\tau y,B(p)]$ は、$p$ によらない貼合せ単元と
$p$ によるスカラーを除き一致する（角括弧は二項一次の行列式で、$C_2$ 側は $\tau$ で座標を運ぶ）。
両者は $p$ の二次式で、どの $p$ においても二つの一次因子の積は長さ4の $T$ 上で恒等零にならない。
従って $p$ によるスカラーは両側で極も零も持たず定数となる。
$T$ 上の $p$ の二つの根を因子として比較すれば $(**)$ であり、非被約な $T$ でも有効である。
$F=A^{-1},G=B^{-1}\tau$ と置く。$T_0$ では $(**)$ の同一の一次因子を消して $F=G$ となる。
局所アフィン座標で二次式に $p=x,y$ を代入すれば $(Fx-x)(y-x)=0$、
$(Gy-y)(y-x)=0$。従って $F,G$ は $T_1$ を各々点ごとにスキームとして固定する
（$y$ による $T_1$ の像も同じ因子である）。
$T_0\cup T_1$ の長さは3以上なので $F=G$ は射影同型として成立し、$B=\tau A$ となる。

ゆえに $z\in C_1$ と $\tau z\in C_2$ を結ぶ一般の直線は三割線である。三番目の点が乗る成分を
$C_1$ と呼び直せる（対合なのでこの直線族は同じである）。その点を $g(z)$ とすれば $g$ は射影同型である。
実際、$C_1$ の三つのHankel二次式を $C_2$ に制限し、長さ4の共通零因子を除くと
基点自由な二次式系が得られる。これが $C_2$ の点を通る $C_1$ の唯一の割線の
二次因子へ送る写像である。Hankel行列の核の二項二次式を取れば、相異なる根の二評価ベクトルを結ぶ直線が戻り、
重根では接線が戻る。$z$ と $g(z)$ の非順序対の系の次数が従って2であるから $1+\deg g=2$ となる。
この議論は三番目の点が $z$ に重なる場合も含む。

この二次因子族 $z+g(z)$ の割線に $\tau z$ が乗る包含は全基底で閉条件として成立する。
$T$ 上で $\tau z$ 自身が $C_1$ にあるので、$z=x$ とすれば $y$ は因子 $x+g(x)$ に入る。
割線上の $C_1$ のスキームは普遍長さ2因子そのものである。実際、二項二次式を $Q$ とすれば
その割線の二平面式は $Qs,Qt$ であり、$\mathbb P^1$ の $s,t$ は基点を持たない。
従って $(y-x)(y-gx)=0$、$T_1$ 上で $gx=y$ である。
$T_1$ が二単点なら $g$ は二点を交換し、二重点なら一点を固定して接倍率が $-1$ である。
どちらも $g$ は非自明な対合となる。
その割線は先に示した滑らかな二次曲面を掃き、$C_2$ もそこに入るので矛盾する。
\end{proof}

\begin{corollary}\label{cor:point-curve-high-coverage}
点ファイバーの一次元点グラフで純部分の次数が4以上なら、全標準AS窓は既出の有限アトラスで被覆される。
六次では直線、二次曲線、平面三次のいずれかの成分が必要である。
全次数を通じ、点グラフ一次元の枝でこれらの印、四次完全交叉印、整五次・七次の既出分類の
どれにも入らないとすれば、純部分は空間有理正規三次で有限余長6を持つ。
\end{corollary}
\begin{proof}
六次は系\ref{cor:all-septic-point-windows}と本定理による。
四次で直線・二次曲線がなければ純部分は整であり、系\ref{cor:integral-hull-high-degrees}を適用する。
五次・七次は既出の全被覆、三次の残余の有限長は定理\ref{thm:pure-curve-defect}である。
\end{proof}
